%% extraction.tex -- the two-word hole problem with target self events, forced
%% target heads (FH), the extracted instances I_pq with enriched rejection
%% (neither dead by support nor root-dead), FH by construction, and the theorem
%% that under satisfiability and (U), condition (II) holds iff no extracted
%% instance has a hole. Proofs: appendix-summaries.tex (app:sm).

\section{Two-Word Holes with Self Events}\label{sec:ex}

Condition \CondII\ fails when a visible atom that is live in $\varphi$ has a
visible $\psi$-consequence outside its closure in $\varphi$; by
\cref{lem:ch-pure}(2), the consequence can be taken to be derived by a pure
derivation. By the lowest point (\cref{lem:sm-lowest-point}), the seed and the
consequence then share a suffix, which follows a pop head in the seed and a
push head in the consequence. Liveness of the seed is read by the source
automaton of \cref{sec:un} together with the comparisons of the seed word
that are self events (\cref{thm:sm-dead}); membership in the closure is read
by the lower automata of $\varphi$ together with comparisons between the two
words (\cref{lem:sm-membership}); and the derivation is read by the bound
automata of $\psi$. This section abstracts these tests into a hole problem
for pairs of words read by four automata (\cref{sec:ex-problem}), extracts its
instances from $\varphi$ and $\psi$ (\cref{sec:ex-instances}), and proves
that, once satisfiability and \CondU\ are known, \CondII\ holds iff no
extracted instance has a hole (\cref{sec:ex-reduction}). The extracted
instances read the seed with a rejection test that already excludes the
root-dead words of $\psi$. With this \emph{enrichment} they have forced target
heads by construction, the property on which \cref{sec:nm} relies. Proofs
omitted here are in \cref{app:sm}.

\subsection{The Hole Problem}\label{sec:ex-problem}

\begin{definition}[Two-word hole problem with self events]\label{def:ex-ss2self}\label{def:ex-ss2}
An \emph{instance} of the \emph{two-word hole problem with self events} is a
tuple
\[
 I=(A,\Bad,B,\Rej,\Guard,\GuardS,A',B',\Head)
\]
that consists of complete deterministic automata $A$, $B$, $A'$, $B'$ over $\Sigma$, with state sets
$S_A$, $S_B$, $S_{A'}$, $S_{B'}$ and a start state $0$ each, together with
sets $\Bad\subseteq S_A$, $\Rej\subseteq S_B$, $\Guard\subseteq S_A\times S_B$,
$\GuardS\subseteq S_B\times S_B$, and $\Head\subseteq S_{A'}\times S_{B'}$.
We call $A$ the \emph{spine automaton}, $B$ the \emph{target automaton}, $A'$
and $B'$ the \emph{head automata}, $\Guard$ the \emph{guard}, $\GuardS$ the
\emph{self guard}, and $\Head$ the \emph{head relation}, and we put
$N_A=\len{S_A}$, and similarly $N_B$, $N_{A'}$, $N_{B'}$. For words $u$ and
$t$, let $A_s(u)$ be the state of $A$ after $\wpre us$ and $B_e(t)$ the state
of $B$ after $\wpre te$, define $A'_s(u)$ and $B'_e(t)$ in the same way, and
write $M(w)$ for the state of an automaton $M$ after the whole word $w$. If
some cut $e$ of $t$ has $B_e(t)\in\Rej$, the least one is the \emph{first
rejected cut} $J(t)$. A pair of words $(u,t)$ is a \emph{hole} of $I$ iff
\begin{enumerate}[label=(\alph*),leftmargin=*]
\item $A_s(u)\notin\Bad$ for every cut $s\le\len u$;
\item $J(t)$ is defined;
\item there are no cuts $s\le\len u$ and $e\le J(t)$ with
  $\wsuf te\pre\wsuf us$ and $(A_s(u),B_e(t))\in\Guard$;
\item there are cuts $s\le\len u$ and $e\le\len t$ with $\wsuf us=\wsuf te$
  and $(A'_s(u),B'_e(t))\in\Head$;
\item there are no cuts $i<j\le J(t)$ with $\wsuf tj\pre\wsuf ti$ and
  $(B_i(t),B_j(t))\in\GuardS$.
\end{enumerate}
The \emph{hole problem} asks whether a given instance has a hole.
\end{definition}

We call $u$ the \emph{spine} and $t$ the \emph{target}. Unlike rejection in
\cref{def:bg-hole}, condition (b) requires a rejected prefix of $t$, not
necessarily a rejected final state; for extracted instances these conditions
coincide, as shown in \cref{sec:ex-instances}. A state of $A$ outside
$\Bad$ is \emph{safe}; a word is \emph{safe from} a state $r$ if $A$ visits
only safe states while reading it from $r$, including $r$ and the last state;
and the spine $u$ is \emph{safe} if it is safe from $0$, which is condition
(a). The other conditions are stated in terms of the following pairs of cuts.

\begin{definition}[Comparisons and events]\label{def:ex-events}
Let $u$ and $t$ be words. A \emph{cross comparison} of $(u,t)$ is a pair of
cuts $(s,e)$ with $0\le s\le\len u$, $0\le e\le\len t$, and
$\wsuf te\pre\wsuf us$. Its \emph{offset} is
$\Delta(s,e)=(\len u-s)-(\len t-e)\ge0$, its \emph{overhang} $o(s,e)$ is the
suffix of $u$ of length $\Delta(s,e)$, so that $\wsuf us=\wsuf te\,o(s,e)$,
and it is \emph{aligned} if $\Delta(s,e)=0$. A \emph{self comparison} of $t$
is a comparison $(i,j)$ of $t$ (\cref{def:bg-comparison}): $i<j\le\len t$ and
$\wsuf tj\pre\wsuf ti$; its overhang is the word $o$ of length $j-i$ with
$\wsuf ti=\wsuf tj\,o$. For an instance $I$ and a target $t$ with $J(t)$
defined:
\begin{itemize}[leftmargin=*]
\item an \emph{event} of $(u,t)$ is a cross comparison $(s,e)$ with
  $e\le J(t)$ and $(A_s(u),B_e(t))\in\Guard$;
\item a \emph{self event} of $t$ is a self comparison $(i,j)$ of $t$ with
  $j\le J(t)$ and $(B_i(t),B_j(t))\in\GuardS$;
\item a \emph{head event} of $(u,t)$ is an aligned cross comparison $(s,e)$
  with $(A'_s(u),B'_e(t))\in\Head$.
\end{itemize}
\end{definition}

Thus a hole is a pair with a safe spine, a target with $J(t)$ defined, no
event, no self event, and a head event. A self comparison $(i,j)$ of $t$ is
the cross comparison $(i,j)$ of the pair $(t,t)$, with offset $j-i>0$; in particular, it
is never aligned. Events relate two words; self events relate two cuts of the
target. An instance with $\GuardS=\varnothing$ has no self events, and its holes
are those of the suffix-sharing two-word problem of a covariant constructor;
\cref{rem:ex-covariant} explains when the extracted instances have this form.

\paragraph{Head form}
A head event $(s,e)$ splits the two words as
\[
 u=\beta\sigma,\qquad t=\gamma\sigma,\qquad
 \beta=\wpre us,\quad\gamma=\wpre te,\quad\sigma=\wsuf us=\wsuf te,
\]
with $(A'(\beta),B'(\gamma))\in\Head$. We call $\beta$ the \emph{spine head},
$\gamma$ the \emph{target head}, and $\sigma$ the \emph{common suffix}. In the
extracted instances of \cref{sec:ex-instances}, the target head is the pop
head of \cref{lem:sm-lowest-point} and the spine head its push head.
\Cref{fig:ex-hole} draws the head form together with the two kinds of pairs
that a hole must avoid. A cross comparison and a self comparison look alike:
in both, a suffix of the target is laid at a cut of a word, the spine or the
target itself, and agrees with that word letter by letter until the suffix
ends.

\begin{figure}[tbp]
\centering
\begin{tikzpicture}[>=Stealth,font=\small,x=0.72cm,y=1cm,
    every node/.style={inner sep=1.5pt},
    box/.style={draw,minimum height=0.5cm,inner sep=0pt},
    lb/.style={font=\footnotesize,inner sep=1pt},
    side/.style={font=\footnotesize,anchor=west,align=left}]
  % spine u = beta sigma and target t = gamma sigma, aligned at their ends
  \draw (0,2.75) rectangle node{$\beta$} (5,3.25);
  \draw[fill=black!12] (5,2.75) rectangle node{$\sigma$} (9,3.25);
  \node[lb,anchor=east] at (-0.15,3) {spine $u$};
  \draw (3,1.75) rectangle node{$\gamma$} (5,2.25);
  \draw[fill=black!12] (5,1.75) rectangle node{$\sigma$} (9,2.25);
  \node[lb,anchor=east] at (2.85,2) {target $t$};
  \node[lb,anchor=south] at (5,3.25) {$s_0$};
  \node[lb,anchor=north east] at (5,1.75) {$e_0$};
  \draw[densely dotted,thick] (5,2.25) -- (5,2.75);
  \draw[densely dotted] (8,2.3) -- (8,1.25);
  \node[lb,anchor=north] at (8,1.25) {$J(t)$};
  \node[side] at (9.3,2.5) {head event $(s_0,e_0)$:\\
    $\wsuf u{s_0}=\wsuf t{e_0}=\sigma$,\\ $(A'_{s_0}(u),B'_{e_0}(t))\in\Head$};
  % a cross comparison (s,e): t[e:] placed at the cut s of u
  \draw[box,dashed] (1,0.25) rectangle node{$\wsuf te$} (6,0.75);
  \draw[densely dotted] (6,0.25) rectangle node{$o$} (9,0.75);
  \draw[densely dotted] (1,2.75) -- (1,0.75);
  \node[lb,anchor=south east] at (1,3.25) {$s$};
  \node[lb,anchor=east] at (0.85,0.5) {$\wsuf te$ at $s$};
  \node[side] at (9.3,0.5) {event if $e\le J(t)$\\ and $(A_s(u),B_e(t))\in\Guard$};
  % a self comparison (i,j): t[j:] placed at the cut i of t
  \draw[box,dashed] (6,-0.75) rectangle node{$\wsuf tj$} (8,-0.25);
  \draw[densely dotted] (8,-0.75) rectangle node{$o$} (9,-0.25);
  \draw[densely dotted] (6,1.75) -- (6,-0.25);
  \node[lb,anchor=north west] at (6,1.75) {$i$};
  \node[lb,anchor=east] at (5.85,-0.5) {$\wsuf tj$ at $i$};
  \node[side] at (9.3,-0.5) {self event if $j\le J(t)$\\ and $(B_i(t),B_j(t))\in\GuardS$};
\end{tikzpicture}
\caption{The pairs of cuts that decide whether $(u,t)$ is a hole
(\cref{def:ex-ss2self,def:ex-events}). Top: the spine and the target, aligned
at their ends; a head event $(s_0,e_0)$ splits them into the heads $\beta$ and
$\gamma$ and the common suffix $\sigma$ (shaded). Middle: a cross comparison
$(s,e)$ places $\wsuf te$ at the cut $s$ of $u$, where it must be a prefix of
$\wsuf us$, leaving the overhang $o$. Bottom: a self comparison $(i,j)$ places
$\wsuf tj$ at the cut $i$ of $t$ in the same way. A hole needs a head event,
a safe spine and a first rejected cut $J(t)$ of the target, and it must have
no event and no self event: every cross comparison $(s,e)$ with $e\le J(t)$
must carry a state pair outside the guard $\Guard$, and every self comparison
$(i,j)$ with $j\le J(t)$ a pair outside the self guard $\GuardS$.}
\Description{Two bars at the top, the spine u and below it the shorter target
t, aligned at the right. The spine consists of beta and the shaded common
suffix sigma, the target of gamma and sigma; a dotted line joins the cuts s
zero and e zero, where sigma begins, and is labeled as the head event. A
dotted line marks the first rejected cut J of t. Below, a dashed box holding
the suffix of t from the cut e starts at the cut s of the spine and is
followed by the overhang o; it is labeled event if e is at most J and the
guard holds. At the bottom, a dashed box holding the suffix of t from the cut
j starts at the cut i of the target and is followed by the overhang o; it is
labeled self event if j is at most J and the self guard holds.}
\label{fig:ex-hole}
\end{figure}


\begin{definition}[Forced target heads]\label{def:ex-fh}
An instance $I$ has \emph{forced target heads}, condition \emph{(FH)}, if for
every $(r,r')\in\Head$ with $r$ reachable in $A'$ and every word $\gamma$ with
$B'(\gamma)=r'$, no proper prefix of $\gamma$ is rejected:
$B_i(\gamma)\notin\Rej$ for all $i<\len\gamma$.
\end{definition}

\begin{lemma}[Heads end before the first rejected cut]\label{lem:ex-head-before-J}
Let $I$ satisfy (FH), and let $(s,e)$ be a head event of $(u,t)$, where $J(t)$
is defined. Then $e\le J(t)$: the target head ends at or before the first
rejected cut, and $\wsuf t{J(t)}$ is a suffix of the common suffix.
\end{lemma}

\begin{proof}
The state $A'_s(u)$ is reachable, and $B'(\wpre te)=B'_e(t)$ is paired with
it by $\Head$. By (FH), $B_i(t)=B_i(\wpre te)\notin\Rej$ for every $i<e$, so the
first rejected cut is at least $e$.
\end{proof}

\Cref{sec:nm} uses \cref{lem:ex-head-before-J} to split a hole at its first
rejected cut. In the extracted instances, (FH) holds by construction
(\cref{lem:ex-fh}).

\subsection{The Extracted Instances}\label{sec:ex-instances}

The instance for a pair of visible signed variables $p,q\in X^\pm$ tests a
pure derivation of $\psi$ from a seed $(p,t)$ to an atom $(q,u)$. The names
follow the readers of \citet{Zhao26}: the \emph{target} is the word whose rejection
is tested, here the seed word $t$, read by the source automaton $\Bsrc_p$ of
\cref{sec:un}; the \emph{spine} is the word $u$ of the derived atom. The
derivation therefore runs from the target to the spine.

\begin{definition}[Extracted instance]\label{def:ex-instance}
For $p,q\in X^\pm$, the \emph{extracted instance}
$I_{pq}=I_{pq}(\varphi,\psi)$ consists of:
\begin{itemize}[leftmargin=*]
\item $A$, the lower automaton of $q$ in $\varphi$, with
  $\Bad=\{Q:Q\cap\Fup\varphi\ne\varnothing\}$;
\item $B=\Bsrc_p$, the source automaton of $p$ (\cref{def:un-source}), with
  \[
   \Rej=\{\zeta:\ \mathrm{dt}(\zeta)\text{ and }\mathrm{du}(\zeta)\text{ both
   fail}\};
  \]
\item $\Guard=\{(Q,(P,b,P')):Q\cap P\ne\varnothing\}$, and $\GuardS$ the self
  guard of \cref{def:un-source};
\item $A'$, the lower automaton of $q$ in $\psi$, and $B'$, the upper
  automaton of $p$ in $\psi$, with $\Head=\{(Q',P'):Q'\cap P'\ne\varnothing\}$.
\end{itemize}
The start states $\downc\varphi q$, the start state of $\Bsrc_p$,
$\downc\psi q$, and $\upc\psi p$ are renamed $0$.
\end{definition}

\begin{table}[tbp]
\caption{The extracted instance $I_{pq}$ for a seed $(p,t)$ and a derived atom
$(q,u)$. The state of $B$ at a cut $e$ of $t$ is
$(P^\varphi_{p,e}(t),[(\bar p,\wpre te)\in\mu_\varphi],P^\psi_{p,e}(t))$, by
\cref{lem:un-source}.}
\label{tab:ex-instance}
\centering
\small
\begin{tabular}{@{}lll@{}}
\toprule
Part & Reads & Meaning\\
\midrule
$A$, $\Bad$ & $Q^\varphi_{q,s}(u)$ on the spine $u$ &
  safe spine: $(q,u)\notin\mu_\varphi$ (\cref{lem:sm-mu})\\
$B$, $\Rej$ & the source automaton on the target $t$ &
  rejected: $\neg\DeadT_\varphi(p,\wpre te)$ and $\neg\DeadU_\psi(p,\wpre te)$\\
$\Guard$ & $Q^\varphi_{q,s}(u)\cap P^\varphi_{p,e}(t)$ &
  event: $(q,u)\in \Reach_\varphi(p,t)$ (\cref{lem:sm-membership})\\
$\GuardS$ & $P^\varphi_{p,i}(t)\cap\overline{P^\varphi_{p,j}(t)}$ &
  self event: $(p,t)$ is dead (\cref{thm:sm-dead})\\
$A'$, $B'$, $\Head$ & $Q^\psi_{q,s}(u)\cap P^\psi_{p,e}(t)$ &
  head event: $(p,t)\to^*_\psi(q,u)$ (\cref{lem:sm-lowest-point})\\
\bottomrule
\end{tabular}
\end{table}

\Cref{tab:ex-instance} summarizes the instance. The automata have
$N_A=2^{2\len X}$, $N_B=2^{4\len X+2\len Z+1}$, and
$N_{A'}=N_{B'}=2^{2(\len X+\len Z)}$ states, and the $4\len X^2$ instances are
computable from $\varphi$ and $\psi$.

Rejection is enriched: a cut of the seed is rejected only if it is neither
dead by support in $\varphi$ nor root-dead in $\psi$. Both properties are
inherited by prefixes (\cref{lem:sm-roots}(4), \cref{lem:sm-dead-states}(4)),
so a rejected cut stays rejected on every extension: $J(t)$ is defined iff $t$
itself is rejected, that is, iff neither $\DeadT_\varphi(p,t)$ nor
$\DeadU_\psi(p,t)$ holds.

\begin{lemma}[Forced target heads]\label{lem:ex-fh}
Every extracted instance $I_{pq}$ satisfies (FH). No hypothesis on $\varphi$
or $\psi$ is needed.
\end{lemma}

\begin{proof}
Let $(Q',P')\in\Head$ and $B'(\gamma)=P'$. Then
$P^\psi_p(\gamma)=P'\ne\varnothing$, and by \cref{lem:sm-head-prefix},
$\DeadU_\psi(p,\wpre\gamma i)$ for every $i<\len\gamma$. By
\cref{lem:un-source}(2), $\mathrm{du}$ holds at every such cut, so
$B_i(\gamma)\notin\Rej$.
\end{proof}

Without the enrichment, that is, with $\Rej$ defined by the failure of
$\mathrm{dt}$ alone, the target head can extend beyond the first rejected cut
(\cref{ex:ex-arrow}).

\subsection{Holes and Condition (II)}\label{sec:ex-reduction}

\begin{theorem}[Holes and condition (II)]\label{thm:ex-reduction}
Let $\varphi$ be satisfiable, let \CondU\ hold, let $p,q\in X^\pm$, and let
$u,t$ be words. Then $(u,t)$ is a hole of $I_{pq}$ iff $(p,t)$ is live in
$\varphi$, $(p,t)\to^*_\psi(q,u)$, and $(q,u)\notin\Cl_\varphi(p,t)$.
Consequently:
\begin{enumerate}[leftmargin=*]
\item condition \CondII\ holds iff no extracted instance $I_{pq}$ has a hole;
\item for every hole $(u,t)$ of $I_{pq}$, the least model $\Cl_\varphi(p,t)$
  of $\varphi$ that contains $(p,t)$ is not the restriction to $X$ of any
  model of $\psi$, so its solution has no extension to a solution of $\psi$.
\end{enumerate}
\end{theorem}

The conditions of \cref{def:ex-ss2self} match the three requirements one by
one. A head event is the lowest point of the pure derivation of $\psi$. A safe
spine and the absence of events say that $(q,u)$ lies neither in $\mu_\varphi$
nor in $\Reach_\varphi(p,t)$ (\cref{lem:sm-membership}). A rejected target and the
absence of self events say that $(p,t)$ is neither dead by support nor has a
self event, that is, that it is live (\cref{thm:sm-dead}). Here \CondU\ is used:
a live seed must also be rejected, which requires that it not be root-dead in
$\psi$. The bounds $e\le J(t)$ and $j\le J(t)$ in \cref{def:ex-ss2self} lose
nothing: the guards require a nonempty set $P^\varphi_{p,e}(t)$, and the pop bound
(\cref{lem:sm-roots}(2)) empties these sets after $J(t)$. The proof is in
\cref{app:sm}.

\begin{corollary}[Entailment and holes]\label{cor:ex-entailment}
Let $\varphi$ and $\psi$ be satisfiable. Then $\varphi\models\Exists\psi$ iff
condition \CondI\ holds, condition \CondU\ holds, and no extracted instance
$I_{pq}$ with $p,q\in X^\pm$ has a hole.
\end{corollary}

\begin{proof}
If \CondU\ fails, the entailment fails by \cref{thm:ch-characterization}.
Otherwise, by \cref{thm:ch-characterization}, the entailment holds iff \CondI\
and \CondII\ hold, and by \cref{thm:ex-reduction}(1), \CondII\ holds iff no
$I_{pq}$ has a hole.
\end{proof}

Conditions \CondI\ and \CondU\ are decidable
(\cref{prop:sm-decide,thm:un-decidable}), and every extracted instance
satisfies (FH) (\cref{lem:ex-fh}). It remains to decide the hole problem for
instances with forced target heads, which \cref{sec:nm,sec:mo} do.

\begin{example}[Condition (U) cannot be dropped]\label{ex:ex-u-needed}
Without \CondU, a failure of \CondII\ need not be a hole. Let $\varphi$ be
empty over $X=\{x,y\}$, and let $\psi=\{x\le y,\ x\le f(a,b)\}$ with $a,b$
hidden and $f$ binary and covariant. The seed $(x^0,\eps)$ is live in
$\varphi$, the copy $(x^0,\eps)\to(y^0,\eps)$ is a pure move of $\psi$, and
$(y^0,\eps)\notin\Cl_\varphi(x^0,\eps)$, so \CondII\ fails at this seed. But
$x^0$ is the root of the upper constructor literal $x^0\le f(a^0,b^0)$ of
$\psi^\pm$, so $\DeadU_\psi(x^0,\eps)$ holds and the target $\eps$ of
$I_{x^0y^0}$ is not rejected: $(\eps,\eps)$ is not a hole. Here \CondU\ fails,
with the countermodel $x=\Top$. The enrichment trades \CondU\ for (FH): with
the rejection test given by the failure of $\mathrm{dt}$ alone, the holes still
characterize the failures of \CondII\ when $\varphi$ is satisfiable, without
assuming \CondU, but (FH) can fail (\cref{ex:ex-arrow}).
\end{example}

\begin{remark}[Covariant and joint-parity instances]\label{rem:ex-covariant}
Informally, if $\varphi$ has a parity assignment (\cref{rem:ch-parity}), in
particular if $f$ is covariant, the self guard never applies to reachable
states of the target automaton (\cref{rem:sm-parity}), and the extracted
instances behave like instances of the suffix-sharing two-word problem
without self events.
\end{remark}

\begin{example}[Shift closure: a hole]\label{ex:ex-running}
Continue \cref{ex:un-running} with $\psi=\{x\le z,\ z\le y,\ z\le f(z,z)\}$ and
$\varphi'_1=\{x\le y,\ y\le g,\ g\le f(g,g)\}$, both satisfiable. Condition
\CondU\ holds, and \CondI\ holds because $\psi$ has no facts.
Consider the seed variable $y^1$ and the derived variable $x^1$. The copy and
constructor literals of $(\varphi'_1)^\pm$ are $x^0\le y^0$, $y^0\le g^0$,
$g^0\le f(g^0,g^0)$, $y^1\le x^1$, $g^1\le y^1$, and $f(g^1,g^1)\le g^1$. In
$I_{y^1x^1}$:
\begin{itemize}[leftmargin=*]
\item $A$ starts at $\downc{\varphi'_1}x^1=\{x^1,y^1,g^1\}$ and moves to
  $\{g^1\}$ on every letter; $\varphi'_1$ has no facts, so every spine
  is safe; the state $\{g^1\}$ lies on a loop at which both letters are safe,
  so it is \emph{branching} (\cref{def:nm-kinds}), the case in which
  \cref{sec:nm} completes every pre-hole directly (\cref{lem:nm-breaker});
\item $P^{\varphi'_1}_{y^1}$ is $\{y^1,x^1\}$ at the empty word and empty
  afterwards, $\mu_{\varphi'_1}=\varnothing$, and $\DeadT^\eps_{\varphi'_1}=
  \{x^0,y^0,g^0\}$; with \cref{ex:un-running}, no cut of any target is dead by
  support or root-dead, so $J(t)=0$ for every target $t$;
\item hence the only possible event is $(0,0)$, a cross comparison iff
  $t\pre u$, and there are no self events, since
  $P^{\varphi'_1}_{y^1,j}(t)=\varnothing$ for $j\ge1$;
\item $Q^\psi_{x^1}(\beta)$ always meets $\upc\psi y^1=\{y^1,z^1,x^1\}$, and
  $P^\psi_{y^1}(\gamma)=\varnothing$ for $\gamma\ne\eps$, so the head events
  are the pairs $(s,0)$ with $\wsuf us=t$.
\end{itemize}
So the holes of $I_{y^1x^1}$ are the pairs $(u,t)$ in which $t$ is a suffix but
not a prefix of $u$, for example $(01,1)$, with head event $(1,0)$. By
\cref{thm:ex-reduction}, $\psi$ derives $(x^1,01)$ from the live seed
$(y^1,1)$, and $(x^1,01)\notin\Cl_{\varphi'_1}(y^1,1)$. In words: if $y$ has a
leaf $\Bot$ at a prefix of $1$, then every $z$ between $x$ and $y$ with
$z\le f(z,z)$ has a leaf $\Bot$ at a prefix of $1$ and hence at a prefix of
$01$, and so does $x$. A finite countermodel is
$x=y=f(f(\Bot,f(\Bot,\Bot)),\Bot)$, whose constructor positions are $\eps$,
$0$, and $01$, with $g$ the complete binary tree whose constructor positions
are the words of length at most $2$. The least model $\Cl_{\varphi'_1}(y^1,1)$
of \cref{thm:ex-reduction}(2) is the solution $x=y=f(\Omega,\Bot)$ and
$g=\Omega$. Since $J(t)=0$, the target head is empty, and (FH) holds trivially
for this instance.
\end{example}

\begin{example}[A target head beyond the cuts dead by support]\label{ex:ex-arrow}
Continue \cref{ex:un-arrow} with $\varphi=\{x\le x\to x\}$ and
$\psi=\{x\le a\to x,\ c\to y\le a\}$, both satisfiable: $x=\Bot$ solves
$\varphi$, and $x=y=\Bot$, $a=c=\Top$ solve $\psi$. Conditions \CondI\ and
\CondU\ hold. Consider the seed variable $x^0$, the derived variable $y^1$,
and the pair $(u,t)=(\eps,01)$ in $I_{x^0y^1}$. Along the target $t=01$, the
upper automata give
\[
 P^\varphi_{x^0}:\ \{x^0\},\ \{x^1\},\ \varnothing,\qquad
 P^\psi_{x^0}:\ \{x^0\},\ \{a^1\},\ \{y^1\}
\]
at the cuts $0,1,2$. With $\DeadT^\eps_\varphi=\{x^0\}$ and
$\DeadU^\eps_\psi=\{x^0,a^1\}$, cut $0$ is dead by support, cut $1$ is not dead
by support but root-dead in $\psi$, and cut $2$ is rejected: $J(01)=2$. The
head event is $(0,2)$, since $Q^\psi_{y^1}(\eps)=\{y^1\}=P^\psi_{x^0}(01)$: the
derivation pops $0$ from $x^0$ to $a^1$ and then $1$ from $a^1$ to $y^1$. The
only cross comparison $(s,e)$ with $e\le2$ is $(0,2)$, whose guard
$\{y^1\}\cap P^\varphi_{x^0}(01)$ is empty; the self comparisons of $01$ ending
at $2$ have an empty self guard, and $(0,1)$ is not a self comparison. The
spine is safe because $\varphi$ has no facts. So $(\eps,01)$ is a
hole: $\psi$ derives $(y^1,\eps)$, that is $y=\Bot$, from the live seed
$(x^0,01)$. A countermodel is $x=(\Bot\to\Bot)\to\Bot$ and $y=\Top$: from
$x\le a\to x$ and $c\to y\le a$, any extension would have
$c\to\Top\le_c a\le_c\Bot\to\Bot$, hence $\Top\le_c\Bot$.

The target head is $\gamma=01$. The rejection test given by the failure of
$\mathrm{dt}$ alone would reject already at cut $1$, inside the head: $(x^0,0)$
is not dead by support, although it is dead by the self event $(0,1)$ of the
word $0$ (\cref{ex:sm-arrow}). So the unenriched instance does not satisfy
(FH). The enrichment rejects only at cut $2$, because $(x^0,0)$ is root-dead
in $\psi$.
\end{example}
